\scanpage{117}
% archive sheet number: 116
Ainsi
\[
 \boxed{\Psi_*:=\Psi_{*P}=S\amalg A\amalg F}\tag*{(12)}
\]
devient un ens ordonné de type spécial, permettant de construire des ens de drapeaux de tt type $H\subset\{0,1,2\}$, $H\ne\varnothing$, \deleted{\unsure{comme|avec}} d'où
\[
 \boxed{\Psi_{\unsure{**|\bullet *}}:=\Psi_{**P}}\tag*{(12 bis)}
\]
s'identifiant à une partie de $\Phi_{**}(X)$, \deleted{et} \inserted{l'ens des} drapeaux de type maximum \deleted{\unsure{sur lequel}} \inserted{$\{0,1,2\}$} \deleted{une opération de $\Gamma$ sur $\Psi_{**P}$,} s'appelle l'ens des repères de \deleted{$\Phi_{**}(X)$, et} \inserted{$P$} \deleted{les conditions ci}
\[
 \mathrm{Rep}(P)=\Psi_{\{0,1,2\}P}\tag*{(13)}
\]
et $\Gamma$ y opère par
\[
 \left\{\begin{aligned}
 \sigma_0\{s,a,f\}&=(s',a,f)&s'&\ne s\\
 \sigma_1(s,a,f)&=(s,a',f)&a'&\ne a\\
 \sigma_2(s,a,f)&=(s,a,f')&f'&\ne f
 \end{aligned}\right.\tag*{(14)}
\]
\deleted{Ainsi} D'ailleurs \unsure{la|les} \unsure{application|applications} de restriction
\[
 \mathrm{Rep}(P)\longrightarrow\Psi_H(P)\tag*{(15)}
\]
se factorise en une application can.
\[
 \mathrm{Rep}(P)/\Gamma_H\longrightarrow\Psi_H(P)\tag*{(16)}
\]
\deleted{\unsure{Ceci fait}} On suppose de plus

$\boxed{17}$ Les applications (15) sont surjectives (donc (16), aussi) --- ou ce qui revient au $=$, \deleted{pour} tt \deleted{\unsure{figure}} \inserted{$s$ est contenu \unsure{dans|plus} \unsure{qq} $f$ (ou $a$)} \deleted{\unsure{couple de figures} \unsure{incidentes} de dim \unsure{non} \unsure{consécutives} (sans exclure les \unsure{faces|figures} \unsure{vides} de dim $-1$ ou de dim $\unsure{2|3}$)} \inserted{et tt $f$ \unsure{contient} un $a$ \unsure{\ldots}} \inserted{\unsure{il y en a une de dim} \unsure{str. int. entre|str. int. côtés}.}
